% 题干转写：完成
% 解析转写：完成
% 图形核对：完成
\section{第37届复赛第1题：高铁空气弹簧劲度系数模型}
\subsection{题干}
高铁空气弹簧由主气室、附加气室及阀门组成，如图\ref{fig:cpho37q01}。平衡时主气室压强、体积分别为 $p_{10}$、$V_{10}$，附加气室容积为 $V_2$，大气压强为 $p_0$。主气室对簧上负载的竖直作用力为 $F=(p_1-p_0)A_e$。在小振幅条件下，主气室体积和有效承载面积随负载相对平衡位置的向下位移 $y$ 线性变化：
\[
\frac{\mathrm dV_1}{\mathrm dy}=-\alpha,\qquad
\frac{\mathrm dA_e}{\mathrm dy}=\beta,\qquad \alpha,\beta>0.
\]
求阀门关闭、阀门打开时，分别在等温与绝热条件下空气弹簧劲度系数 $K=\mathrm dF/\mathrm dy$ 与 $A_e$ 的关系。
\begin{figure}[htbp]\centering
\begin{tikzpicture}[scale=0.95,>=Stealth,font=\small]
  \draw[fill=gray!15,thick] (0.7,2.55) rectangle (3.85,3.45);
  \node at (2.28,3.00) {簧上负载};
  \draw[fill=gray!25,thick] (2.05,2.12) rectangle (2.55,2.55);
  \draw[fill=white,thick,rounded corners=8pt] (1.35,0.55) rectangle (3.25,2.12);
  \draw[fill=gray!25,thick] (1.18,0.26) rectangle (3.42,0.55);
  \node at (2.30,1.33) {主气室};
  \draw[fill=white,thick,rounded corners=8pt] (-2.55,0.50) rectangle (-0.78,1.88);
  \draw[fill=gray!25,thick] (-2.70,0.22) rectangle (-0.63,0.50);
  \node at (-1.67,1.17) {附加气室};
  \draw[very thick] (-1.67,1.88) .. controls (-1.50,2.48) and (0.25,2.40) .. (1.35,1.92);
  \draw[fill=white,thick] (-2.10,2.14) rectangle (-1.50,2.36);
  \draw[thick] (-1.80,2.36) -- (-1.80,2.62);
  \node[left] at (-1.90,2.55) {阀门};
  \draw[thick] (-3.1,0.18) -- (3.85,0.18);
  \foreach \x in {-2.95,-2.55,...,3.65} {\draw (\x,0.18) -- +(0.18,-0.12);}
  \draw[<->] (3.78,0.55) -- (3.78,2.12) node[midway,right] {$y$（向下为正）};
\end{tikzpicture}

\caption{第37届复赛第1题高铁空气弹簧结构的 TikZ 重绘。}\label{fig:cpho37q01}
\end{figure}

\subsection{参考解答整理}
\paragraph{统一推导。} 用多方过程方程
\[
p_1V_1^n=\text{常量}
\]
统一处理两种热力学过程。由 $F=(p_1-p_0)A_e$，有
\[
K=(p_1-p_0)\frac{\mathrm dA_e}{\mathrm dy}
+\frac{\mathrm dp_1}{\mathrm dy}A_e.
\]
对过程方程微分，并代入 $\mathrm dV_1/\mathrm dy=-\alpha$、$\mathrm dA_e/\mathrm dy=\beta$，阀门关闭时得到通式
\[
K=\frac{\alpha n p_{10}V_{10}^n}{V_1^{n+1}}A_e
+\beta\left[p_{10}\left(\frac{V_{10}}{V_1}\right)^n-p_0\right].
\]

\paragraph{（1）阀门关闭。} 上下乘客变化缓慢时，气体近似等温，$n=1$：
\[
\boxed{K=\frac{\alpha p_{10}V_{10}}{V_1^2}A_e
+\beta\left(\frac{p_{10}V_{10}}{V_1}-p_0\right).}
\]
剧烈振动时可近似绝热。由 $C_V=\frac52R$，
\[
\gamma=\frac{C_P}{C_V}=\frac{C_V+R}{C_V}=\frac75.
\]
于是
\[
\boxed{K=\frac{7\alpha p_{10}V_{10}^{7/5}}{5V_1^{12/5}}A_e
+\beta\left[p_{10}\left(\frac{V_{10}}{V_1}\right)^{7/5}-p_0\right].}
\]

\paragraph{（2）阀门打开。} 低频、缓慢变形时，主气室与附加气室压强相同，二者作为整体作等温变化：
\[
\boxed{K=\frac{\alpha p_{10}(V_{10}+V_2)}{(V_1+V_2)^2}A_e
+\beta\left[\frac{p_{10}(V_{10}+V_2)}{V_1+V_2}-p_0\right].}
\]
高频、快速变形时，来不及通过连通管流动的气体仅限于主气室，因此仍为主气室的绝热结果：
\[
\boxed{K=\frac{7\alpha p_{10}V_{10}^{7/5}}{5V_1^{12/5}}A_e
+\beta\left[p_{10}\left(\frac{V_{10}}{V_1}\right)^{7/5}-p_0\right].}
\]

\subsection{转写依据}
题干、答案公式和图形分别依据 \texttt{CPHO37-Q01.md}、\texttt{CPHO37-Q01\_formula\_check\_v1.csv} 与题图核对入口。评分分布为第（1）问32分、第（2）问8分。
